\thispagestyle{empty}
\makeatletter
\begin{center}
{\@titlehead}\\[3.5cm]
{\Huge\bfseries Grundlagen der Informatik}\\[6mm]
{\LARGE Kompaktskript}\\[4mm]
{\large Erklärungen, Beispiele und Aufgaben\\zu Vorlesung, Übungen und Klausur}\\[2cm]
\begin{tikzpicture}[syntax]
	\node[knoten] (a) at (0,0) {\itshape g};
	\node[knoten] (b) at (2,0) {\itshape d};
	\node[knoten] (c) at (4,0) {\itshape i};
	\draw[linie, ->] (-1.5,0) -- (a);
	\draw[linie, ->] (a) -- (b);
	\draw[linie, ->] (b) -- (c);
	\draw[linie, ->] (c) -- (5.5,0);
	\draw[linie, ->] (c.south) -- (4,-1) -- (0,-1) -- (a.south);
\end{tikzpicture}
\end{center}
\makeatother
\vfill

\begin{merke}[So ist das Skript aufgebaut]
Jedes Kapitel besteht aus drei Teilen:
\begin{itemize}[itemsep=0.1em]
\item \textbf{Worum geht es?} Das Thema in wenigen Sätzen, das Wichtigste steht im Merkkasten.
\item \textbf{Beispiele} zeigen das Vorgehen Schritt für Schritt.
\item \textbf{Aufgaben} zum Selbstrechnen. Sie orientieren sich an den Übungsblättern und an den Probeklausuren.
\end{itemize}
Aufgaben mit der Markierung \K\ entsprechen in Art und Umfang einer Klausuraufgabe. Aufgaben mit \S\ gehen etwas über den Standard hinaus und sind für alle gedacht, die schon Vorkenntnisse mitbringen.
\end{merke}
\newpage
\tableofcontents
\newpage
